\documentclass[12pt, letterpaper]{report}
\usepackage{import}
\import{C:/Users/83tal/Documents/School}{header.tex}
\graphicspath{C:/Users/83tal/Documents/School/Newtonian Mechanics (PHY-211)/Figures/}
\newcommand{\ihat}{\boldsymbol{\hat{\textbf{\i}}}}
\newcommand{\jhat}{\boldsymbol{\hat{\textbf{\j}}}}
\newcommand{\khat}{\boldsymbol{\hat{\textbf{k}}}}
\newcommand{\R}{\mathbb{R}} %Real numbers
\begin{document}
There is a theorem that extends this idea into the decomposition of vector fields known as \textbf{Helmholtz} \textbf{Decomposition}, also \textbf{Helmholtz' Theorem} or the \textbf{Fundamental} \textbf{Theorem} \textbf{of} \textbf{Vector} \textbf{Calculus}. The full statement and proof are beyond this class, but the idea can be stated as follows:\\
Let \(\mathbf{F} \) be a vector field defined on a domain bounded by a closed surface \(S\subset \mathbb{R} ^3\). \(\mathbf{F} \) can then be expressed as 
\[
    \mathbf{F} = \mathbf{CF} +\mathbf{DF} 
\]
where \(\mathbf{CF} \) is irrotational and \(\mathbf{DF} \) is incompressible.\\
LOL ``BEYOND THE SCOPE OF THIS CLASS''
\section{Attempt to Understand Helmholtz Decomposition}
\begin{theorem}[Helmholtz Decomposition]
    Let \(V\subseteq \mathbb{R} ^3\) have a boundary \(\partial V=S\). Let \(\mathbf{F}:V\to \mathbb{R} ^3 \) be a vector field with continuous second derivatives on \(V\). \(F\) can then be decomposed into curl-free and divergence-free components as:
    \[
        \mathbf{F} =-\nabla \Phi +\nabla \times \mathbf{A} 
    \]
    where
    \[
        \Phi (\mathbf{r} )=\frac{1}{4\pi }\int_V \frac{\nabla ^{\prime} \cdot \mathbf{F} \left( \mathbf{r} ^{\prime}  \right) }{\left\vert \mathbf{r} -\mathbf{r} ^{\prime}  \right\vert }\,dV^{\prime} -\frac{1}{4\pi }\oint_S \mathbf{\hat{n}}^{\prime} \cdot \frac{\mathbf{F} \left( \mathbf{r} ^{\prime}  \right) }{\left\vert \mathbf{r} -\mathbf{r} ^{\prime}  \right\vert }\,dS^{\prime}
    \]
    \[
        \mathbf{A}(\mathbf{r} )=\frac{1}{4\pi }\int _{V}\frac{\nabla ^{\prime} \times \mathbf{F} \left( \mathbf{r} ^{\prime}  \right) }{\left\vert \mathbf{r} -\mathbf{r} ^{\prime}  \right\vert }\,dV^{\prime} -\frac{1}{4\pi }\oint_S \mathbf{\hat{n}}^{\prime} \times \frac{\mathbf{F} \left( \mathbf{r} ^{\prime}  \right) }{\left\vert \mathbf{r} -\mathbf{r} ^{\prime}  \right\vert }\,dS^{\prime} 
    \]
\end{theorem}
This is complicated so consider the following definitions.
\subsection{Lebesgue Integration}
\begin{definition}[\(\sigma \)-algebra]
    A \(\sigma \)-algebra on some set \(X\) is a nonempty set \(\Sigma \) of subsets of \(X\) that is closed under compliment, countable unions, and countable intersections.
\end{definition}	
An example is for the set \(X=\{ a,b,c,d \} \), the set \(\Sigma =\{ \varnothing ,\{ a,b \},\{ c,d \},\{ a,b,c,d \}    \} \) is a \(\sigma \)-algebra. In general, all finite algebras are \(\sigma \)-algebras.
\begin{definition}[Set Function]
    Let \(\Omega \) be a set. If \(\mathcal{F} \) is a family of sets, meaning \(\mathcal{F} \subseteq \mathcal{P} (\Omega )\), then a function \(\mu :\mathcal{F} \to \mathbb{R} \cup \{ -\infty,\infty \} \) is a set function on \(\mathcal{F} \).
\end{definition} 
\begin{definition}[Measure]
    Let \(X\) be a set with \(\sigma \)-algebra \(\Sigma \). A set function \(\mu :\Sigma \to \mathbb{R} \cup \{ -\infty,\infty \} \) is called a measure if the following conditions hold.
    \begin{itemize}
        \item For all \(E\in \Sigma \), \(\mu (E)\geq 0\).
        \item \(\mu (\varnothing )=0\).
        \item For all countable collections \(\{ E \}_{k=1}^\infty \) of pairwise disjoint sets in \(\Sigma \),
        \[
            \mu \left( \biguplus_{k=1}^{\infty}E_k  \right) =\sum_{k=1}^\infty \mu \left( E_k \right)  
        \]
    \end{itemize}
    where pairwise disjoint sets are sets \(E,F\) such that \(E\cap F=\varnothing \) always, and \(E\uplus F\) denotes the union of disjoint sets.
\end{definition}
\begin{definition}[Mass of a Set Function]
    Let \(\mu \) be a set function on a set \(\mathcal{F} \subseteq \mathcal{P} (S)\) where \(\mathcal{F} \) is a family of subsets of \(S \). Define
    \[
        \cup \mathcal{F} \coloneqq \bigcup_{F\in\mathcal{F} } F 
    \]
    and note \(\cup \mathcal{F} \in \mathcal{F} \). The mass of \(\mu \) is defined as \(\mu (\cup \mathcal{F} )\).
\end{definition}
\begin{definition}[Measurable Functions]
    Let \(E\) be a set with \(\sigma \)-algeba \(X\) and measure \(\mu \) on \(E\). A real-valued function \(f\) on \(E\) is measurable if the preimages of every interval of the form \((t,\infty)\) is in \(X\):
    \[
        \left\{ x \mid f(x)>t \right\} \in X\,\forall t\in\mathbb{R} 
    \]
    For Euclidean spaces, we generally say \(\mu \) is the Lebesgue Measure.
\end{definition}
\begin{definition}[Infimium]
    Let \(S\) be a subset of a poset \(P,\leq \). Some \(a\in P\) is a lower bound of \(S\) if
    \[
        a \leq x\,\forall x\in S
    \]
    Note that it is not required that \(a\in S\). An element \(a\) is called an infimium of \(S\) if
    \[
        \text{For all lower bounds } y\text{ of } S\text{ in } P,\, y\leq a
    \]
    and we write \(\inf S=a\). This can be interpreted as the greatest lower bound of \(S\).
\end{definition}	
\begin{definition}[Lebesgue Outer Measure]
    Let \(I=[a,b ]\) or \(I=(a,b)\) be an interval of the real numbers. Let \(\ell (I)=b-a\) define the length of \(I\). For any subset \(E \subseteq \mathbb{R} \), the Lebesgue Outer Measure \(\lambda^*(E)\) is defined as 
    \[
        \lambda^*(E)=\inf \left\{ \sum_{k=1}^\infty \ell \left( I_k \right) \mid \left( I_k \right)_{k\in\mathbb{N} }\text{ is a sequence of open intervals with }E\subset \bigcup_{k=1}^{\infty}I_k     \right\} 
    \]
\end{definition}
The above definition generalizes to higher dimensions. Let \(C=I_1 \times \cdots I_n\) be the cuboid formed from the cartesian product of open intervals, and let \(\operatorname{vol}(C)=\ell (I_1)\cdots \ell (I_n) \) denote the volume of the cuboid. In this case, for some \(E \subseteq \mathbb{R} ^n\),
\[
    \lambda^*(E)=\inf \left\{ \sum_{k=1}^\infty  \operatorname{vol}\left( I_k \right) \mid \left( C_k \right)_{k\in\mathbb{N} }\text{ is a sequence of products of open intervals with }E\subset \bigcup_{k=1}^{\infty}C_k     \right\} 
\]
The set of intervals/cuboids forms a \(\sigma \)-algebra.
\begin{remark}
    For \((a,b),[a,b]\subset \mathbb{R} \) as intervals of \(\mathbb{R} \), \(\inf (a,b)=\inf [a,b]=a\).
\end{remark}
\begin{definition}[Lebesgue Measure]
    Let \(\lambda^*\) denote the Lebesgue Outer Measure. A set \(E \subseteq \mathbb{R}^n \) is Lebesgue Measurable if
    \[
        \lambda ^*(A) = \lambda ^*(A\cap E)+\lambda ^*\left(A\cap E^{\complement} \right)
    \]
    for any \(A \subseteq \mathbb{R} ^n\). This criteria is called the Carath\'eodory's criterion. The set of all \(E\) satisfying this criterion are the set of all Lebesgue measurable sets. The Lebesgue measure is defined as the outer measure \(\lambda (E)=\lambda ^*(E)\). The set of all \(E\) forms a \(\sigma \)-algebra.
\end{definition}
\begin{intuition}
    In the Lebesgue outer measure, \(E\) is ``covered'' by the sets \(I_k\), since \(E\) is a subset of their union. So the sum of the lengths of these sets must be at least the ``length'' of \(E\), but since \(E\) is only a subset of the union of \(I_k\), there could be points outside of \(E\) being considered, and thus the union may overestimate the length of \(E\). To fix this, a set is constructed of all these lengths and the infimium is taken, which takes the lowest sum of lengths of the unions of intervals that at least covers all of \(E\). This can be thought of as the sum of the lengths of every single individual disjoint nonoverlapping set that composes \(E\).\\
    The second condition prevents strange sets with curious properties from being Lebesgue measurable. For some set \(A\), if the length of the part of \(A\) that intersects with \(E\) and the part that does not intersect with \(E\) sum to the total length of \(A\), then \(E\) is said to be Lebesgue-measurable.
\end{intuition}
We can now begin to define Lebesgue integration. Examples will mostly focus on \(\mathbb{R} \), but results are generalizable. Let \(S\) be a set with measure \(\mu \), such as \(\mathbb{R} \) and the Lebesgue measure.
\begin{definition}[Integral of the Indicator Function]\label{lebint1}
    The integral of the identity function \(1_S\) is defined as 
    \[
        \int 1_S\,d \mu =\mu (S)
    \]
\end{definition}
Since the indicator function is defined as 
\[
    1_S(x) =\left\{
    \begin{array}{cc}
    1&x\in S\\
    0&x\notin S
\end{array}
    \right.
\]
it makes sense to give such a definition, since if we think of the measure as the length, then the area under the curve of a function with height \(1\) on some domain is simply the length of the domain.
\begin{exercise}
    \[
        \int_{[0,1]}1_\mathbb{Q}\,d \mu 
    \]
\end{exercise}
\begin{solution}
    Consider \(1_\mathbb{Q} \), the Dirichlet function. This is Lebesgue integrable from \ref{lebint1}. For \(1_\mathbb{Q} :[0,1]\to \mathbb{R} \), we say
    \[
        \int_{[0,1]} 1_\mathbb{Q} \,d \mu =\mu ([0,1]\cap \mathbb{Q} )
    \]
    We must now find the Lebesgue measure of \(\mathbb{Q} \) over the interval \([0,1]\subset \mathbb{R} \). For the purposes of proper indexing, let \(\phi :\mathbb{N}\to \mathbb{Q}\) be a bijection and let
    \[
        E_k = \left\{ x \mid 0<\vert x-\phi (k) \vert<\varepsilon   \right\} 
    \]
    be the \(\varepsilon \)-neighborhood of \(\phi (k)\). Such a bijection exists because \(\vert \mathbb{N} \vert =\vert \mathbb{Q} \vert=\aleph_0  \)  Since \(\phi \) is a bijection, this gives an \(\varepsilon \)-neighborhood for every rational number. Since the measure is continuous under continuous transformations of \(\varepsilon \),
    \[
        \lim_{\varepsilon  \to 0} \mu (E_k)=0
    \]
    As such, we can say that if we let \(\varepsilon \to 0\),
    \[
        \biguplus_{k=1}^\infty E_k = \mathbb{Q}
    \]
    and
    \[\ell (E_k)=0\,\forall k\]
    Therefore we have 
    \[
        \mu ([0,1]\cap \mathbb{Q})=0
    \]
    \[
        \therefore \int_{[0,1]}1_\mathbb{Q}\,d \mu =0
    \]
    This feels fantastically nonrigorous, hopefully its at least correct tho lol
\end{solution}
\begin{definition}[Simple Functions]
    A finite linear combination of indicator functions 
\[
    \sum_{k}a_k 1_{S_k}
\]
with coefficients \(a_k\) and disjoint measurable sets \(S_k\) is called a measurable \textbf{simple} \textbf{function}. Linearity extends the integral to \textbf{nonnegative} measurable simple functions with coefficients \(a_k\) all positive. We define
\[
    \int\left( \sum_{k}a_k 1_{S_k}  \right) =\sum_{k}a_k \int1_{S_k} \,d \mu =\sum_{k}a_k \mu \left( S_k \right)
\]
and this sum will be some element of the nonnegative extended real number line.
\end{definition}
In the above definition, we assume \(\mu \left( S_k \right)<\infty \) for all \(a_k \neq 0\) to make a definition
\[
    f=\sum_{k}a_k 1_{S_k} 
\]
make sense for an integral.
\begin{definition}
    If \(B\) is a measurable subset of \(E\) and \(s\) is a measurable simple function
    \[
        s=\sum_{k}a_k 1_{S_k} 
    \]
    one defines
    \[
        \int_B s\,d \mu =\int 1_B s\,d \mu =\sum_{k}a_k \mu \left( S_k \cap B \right)
    \]
\end{definition}
Now we can consider non-negative and non-simple functions.
\begin{definition}[Non-negative functions]
    Let \(f\) be a non-negative measurable function on \(E\) which takes values on the extended real number line. We define
    \[
        \int_E f\,d \mu =\sup \left\{ \int_E s\;d \mu \mid 0\leq s\leq f,s\text{ is simple}  \right\}
    \]
\end{definition}
It can be shown that this definition is equivalent to the above definition. For some function \(f\), let \(s_n\) be the simple function with value \(\frac{k}{2^n}\) where
\[
    \frac{k}{2^n}\leq f(x)<\frac{k+1}{2^n}
\]
for a nonnegative integer \(k\). It's fairly obvious that 
\[
    \int f\,d \mu =\lim_{n \to \infty} \int s_n\,d \mu 
\]
This approach of defining a linear combination of indicator functions with attached coefficients is directly analogous to a Riemann sum, but the abstraction puts less restrictions over what functions it can be defined.
\begin{definition}[Signed Functions]\label{lebsign}
    Let \(f\) be a measurable function on \(E\) to the extended reals. Define 
    \[
        f^{+}(x)=\left\{
            \begin{array}{cc}
                f(x)&f(x)>0\\
                0&f(x)\ngtr 0
            \end{array}
        \right. 
    \]
    \[
        f^{-}(x)=\left\{
            \begin{array}{cc}
                -f(x)&f(x)<0\\
                0&f(x)\nless 0
            \end{array}
        \right. 
    \]
    It follows that \(f=f^+ -f^-\). We say that the lebesgue integral exists if at least one of the integrals is finite:
    \[
        \min \left( \int f^+\,d \mu ,\int f^- \,d \mu  \right) <\infty
    \]
    If this is true, we define 
    \[
        \int f\,d \mu =\int f^+ \,d \mu -\int f^- \,d \mu 
    \]
\end{definition}
This definition comes from the fact that the measure is defined to be positive, but we wish to preserve negative values of the integral. This definition gives the desired properties of the integral.\\
We can also say that in the case of some measurable and non-negative \(f\), the function 
\[
    f^* (t)\coloneqq \mu \left( \left\{ x\in E \mid f(x)>t \right\}  \right) 
\]
is monotonic, and the lebesgue integral is given as 
\[
    \int_E f\,d \mu \coloneqq \int_0^\infty f^*(t)\,dt
\]
via the improper Riemann integral. A Lebesgue integrable function that is not necessarily nonnegative is given by the decomposition seen in \ref{lebsign}.\\
\[
    \phi :\Xi \to \mathbb{Z}/8\mathbb{Z}
\]
\[
    \psi \in \Xi 
\]
\[
    \psi \mapsto 1
\]
\[
    \langle \psi  \rangle =\Xi 
\]
\[
    \mu :\mathcal{P} (\Xi )\to \mathbb{R}\cup \left\langle -\infty,\infty \right\rangle 
\]
\[
    \mu (\Omega )=\sup \left\{ \phi (k) \mid k\in \Omega  \right\}
\]
\[
    \Omega \in \mathcal{P} (\Xi)
\]
\[
    \mu (\varnothing )=0
\]
\[
    \sigma_k \coloneqq \nu  ^{-1}  (k)\in \Xi
\]
\[
    k\in \mathbb{Z}/8\mathbb{Z}
\]
\[
    \eta \circ \phi = \phi \circ \eta  =\text{id} 
\]
\[
    1:\Omega \to \Xi 
\]
\[
    \int 1\,d \mu = \mu (\Omega )
\]
\[
    \int 1\circ 1^{\prime} \,d \mu \coloneqq \int1\,d \mu \circ \int1\,d \mu 
\]
\[
    a\in \mathbb{Z}
\]
\[
    \zeta :\Xi \to \Xi 
\]
\[
    \zeta \left( \sigma (k) \right) =2\sigma (k)+ k \sigma (3)+\psi 
\]
\[
    a \sigma (k) = \sigma (k) \cdots \sigma (k)
\]
\[
    \lambda : \Xi \to \mathcal{P} (\Xi )
\]
\[
    \lambda (\sigma (k))=\inf \left\{ \Lambda \in \mathcal{P} (\Xi ) \mid \Lambda \cap \left\{ \left\{ \sigma (k) \right\}\cup \left\{ \psi  \right\}   \right\} = \left\{ \psi ,\sigma (k) \right\}    \right\}
\]
\[
    \Omega \leq  \Omega^{\prime} \iff |\Omega | \leq  \vert \Omega ^{\prime}  \vert 
\]
\end{document}